% ECONOMICS_PROBLEM: {"id": "D47-174", "title": "Scaling kidney exchange", "jel_code": "D47", "source_id": "OP-0298"}
% !TeX program = xelatex
\documentclass[11pt,a4paper]{article}
\usepackage[margin=23mm]{geometry}
\usepackage{fontspec}
\setmainfont{Latin Modern Roman}
\usepackage{amsmath,amssymb,mathtools}
\usepackage{xcolor,enumitem,booktabs,longtable,array,xurl,hyperref}
\definecolor{muted}{HTML}{53636C}
\hypersetup{colorlinks=true,urlcolor=blue,linkcolor=blue}
\setlength{\parindent}{0pt}
\setlength{\parskip}{5pt}
\newcommand{\R}{\mathbb R}
\newcommand{\E}{\mathbb E}
\newcommand{\N}{\mathbb N}
\newcommand{\ind}{\mathbf 1}
\newcommand{\Var}{\operatorname{Var}}
\newcommand{\Cov}{\operatorname{Cov}}
\newcommand{\supp}{\operatorname{supp}}
\DeclareMathOperator*{\argmin}{arg\,min}
\DeclareMathOperator*{\argmax}{arg\,max}
\newcommand{\entrymeta}[3]{{\sffamily\small\color{muted}\textbf{\texttt{#1}}\quad|\quad #2\par #3\par}}
\newcommand{\Problemref}[1]{\texttt{#1}}
\newcommand{\JELref}[2]{\texttt{#2} (JEL \texttt{#1})}
\begin{document}
\section*{Scaling kidney exchange}
\textbf{Problem ID:} \texttt{D47-174}\quad \textbf{JEL:} \texttt{D47}\par
% BEGIN ECONOMICS STATEMENT
\label{problem:OP-0298}
{\sffamily\small\color{muted}Original subject group: Mechanism design and market design\par}
\label{definitions:problem:30007580}
\textbf{Audit classification:} Published research agenda. Roth 2010 p.4 identifies center participation, and Ashlagi–Roth §6 item 2 explicitly asks about international collaboration. The weighted withholding mechanism is editorial.
\subsection*{Definitions and precise research target}
\textbf{Model and research target.} Let countries or transplant centers \(h=1,\ldots,H\) own sets \(V_h\) of incompatible patient–donor pairs. A directed compatibility graph \(G=(V,E)\), with \(V=\bigcup_hV_h\), has edge \(i\to j\) when donor \(i\) can medically donate to patient \(j\). Feasible exchanges are disjoint cycles of length at most \(K\), together with permitted chains from nondirected donors, satisfying explicit cross-border shipping and surgery-capacity constraints. Each center may withhold pairs and conduct internal exchanges. Let \(U_h(M)\) be its expected number of transplanted patients under mechanism \(M\), and \(U_h^{\mathrm{outside}}\) its attainable internal alternative. Fix nonnegative patient weights \(w_i\) representing an openly stated equity objective. Determine mechanisms maximizing
\[
 E\left[\sum_{i\in V}w_i\mathbf1\{i\text{ receives a transplant}\}\right]
\]
while ensuring \(U_h(M)\geq U_h^{\mathrm{outside}}\) and discouraging profitable pair withholding, under a specified stochastic compatibility and arrival model. Quantify separately gains from larger scale and gains from differences in biological compatibility distributions. Include exchange failures and the absence of unrestricted monetary transfers. Characterize the efficiency or equity loss needed for sustainable institutional participation and test mechanisms against realistic pool composition.

\emph{Definitions and maintained assumptions (editorial unless a source definition is named).} Pair vertices are partitioned by owning centers. Add nondirected-donor vertices with no paired recipient and terminal-recipient vertices where chains may end. A feasible plan is a vertex-disjoint collection of directed cycles of length at most \(K\) and directed chains with declared length, timing and surgery/shipping limits; failure realizations alter which transplants actually occur. Specify stochastic arrivals, compatibility and failure distributions and the information each center observes before submitting a subset \(S_h\subseteq V_h\). A mechanism maps submissions to randomized feasible plans. Full-participation dominance means \(\mathbb E[U_h(V_h,S_{-h})]\ge\mathbb E[U_h(S_h,S_{-h})+U_h^{\rm withheld}(V_h\setminus S_h)]\) for every center, subset and other submissions, under the same permitted internal-exchange timing. Interim incentive compatibility averages instead over other centers’ information. An internal outside option must be computed under identical medical and resource constraints. Patient weights are announced equity primitives; centers’ transplant-count objectives need not share them. Size and biological-heterogeneity effects require separate counterfactual pool constructions.
\subsection*{Variants and assumption changes}
\begin{enumerate}
\item \textbf{One-shot center incentives (editorial).} Fix graph distributions and cycle cap, prohibit money and require dominant full submission. Determine the attainable weighted-efficiency loss relative to a centralized complete-information benchmark.
\item \textbf{Dynamic international pool (source direction, editorial model).} Allow arrivals, exchange failures and center exit with continuation outside options. Find participation-sustaining collaboration rules and compare matched-size pools to isolate biological heterogeneity gains.
\end{enumerate}
\subsection*{References and source audit}
\begin{enumerate}
\item Roth §1 p.4; Ashlagi–Roth §6 item2 p.5467. Supports the stated research area; exact displayed specification is editorial. \url{https://web.stanford.edu/~alroth/papers/NSF%20Grand%20Challenge.SBE%202020.%20Market%20Design.pdf}
\item Roth §1 p.4; Ashlagi–Roth §6 item2 p.5467. Supports the stated research area; exact displayed specification is editorial. \url{https://pubsonline.informs.org/doi/pdf/10.1287/mnsc.2020.3954}
\end{enumerate}
\begin{enumerate}\raggedright
\item\hypertarget{definitions:source:30007580:1}{} Alvin E. Roth. 2010. \emph{Market Design: Understanding markets well enough to fix them when they’re broken}. §1 Fundamental questions, pp.3–4.
\url{https://web.stanford.edu/~alroth/papers/NSF%20Grand%20Challenge.SBE%202020.%20Market%20Design.pdf}.
\item\hypertarget{definitions:source:30007580:2}{} Itai Ashlagi; Alvin E. Roth. 2021. \emph{Kidney Exchange: An Operations Perspective}. §6 Open Questions and Research Directions, item2 International Exchanges, p.5467.
\url{https://pubsonline.informs.org/doi/pdf/10.1287/mnsc.2020.3954}.
DOI: \url{https://doi.org/10.1287/mnsc.2020.3954}.
\end{enumerate}

\subsection*{Common conventions: Source mathematical conventions}

These conventions apply to each entry unless explicitly replaced.

\textbf{Models and identification.} A primitive model \(m\) specifies all probability laws, feasible choices, information, equilibrium and measurement rules used by its target. The admissible class \(\mathcal M\) is fixed independently of outcomes. If \(P_O(m)\) is its observable probability law and \(\tau(m)\) the target, its identified set at observable law \(P\) is
\[
 \mathcal I_\tau(P)=\{\tau(m):m\in\mathcal M,\ P_O(m)=P\}.
\]
Point identification means that this set is a singleton. An empty set signals incompatibility of data and assumptions. Coordinate bounds need not describe the joint set. Bounds are sharp only if every retained target is attainable or arbitrarily approachable by compatible models. Finite bounds require explicit restrictions on unobserved quantities; unrestricted confounding, hidden wealth, extrapolation or arbitrary equilibrium selection can yield uninformative sets.

\textbf{Causal interventions.} A potential outcome \(Y_i(p)\) is the outcome under a complete policy regime \(p\), including timing, financing, information and market spillovers. The target population and units are fixed before treatment. With interference, use the complete assignment vector or a justified exposure mapping; individual no-interference notation alone cannot identify an economy-wide intervention. A difference of mean potential outcomes does not require the cross-world joint distribution. Conditioning on post-treatment migration, survival, enrollment or employment changes the estimand and needs separate justification.

\textbf{Statistical statements.} An estimator, confidence set, forecast or test specifies a sampling law, model class and loss/coverage criterion. Uniform \((1-\alpha)\)-coverage means \(\inf_{m\in\mathcal M}\Pr_m\{\tau(m)\in C_n\}\geq1-\alpha\), or an explicitly asymptotic version. The whole parameter space trivially has coverage, so informativeness requires a stated width, risk or power objective. A forecasting improvement is relative to a named benchmark, information set and evaluation population.

\textbf{Equilibrium and optimization.} An equilibrium comprises feasible plans, optimality, beliefs where needed, budget constraints and all market/resource clearing. The primitives or a stated benchmark define each correspondence \(\mathcal E(p;m)\). Nonemptiness is assumed or proved, never inferred just from compact policy sets. A selected-equilibrium target uses a declared selection rule; a robust target quantifies over all permitted equilibria. Use suprema or infima unless attainment is established. An \(\varepsilon\)-optimal policy is feasible and within \(\varepsilon\) of the finite optimum value. Report an empty feasible set separately.

\textbf{Welfare and accounting.} Normative utility, interpersonal normalization, social weights, numeraire, discounting and the use of public receipts are primitives. Behavioral utility need not coincide with the welfare criterion. Household welfare includes ownership income and rents once; adding profits again to compensating variation generally double-counts them. A monetary equivalent is defined by an explicit transfer and price convention, with monotonicity and range conditions. Average effects, mechanism contributions, transfers and welfare losses are different targets. Infinite sums require convergence and dynamic feasible plans require terminal or no-Ponzi conditions.

\textbf{Source versions.} A working-paper date, revision date and journal publication year are distinguished. A DOI for a journal version is not automatically the DOI of an earlier manuscript. Locators refer to the stated version and printed pages where specified. A metadata-only check does not verify a claim about a full-text conclusion. Every entry's source subsection narrows the provenance claim accordingly.
% END ECONOMICS STATEMENT
\end{document}
